%%%%%%%%%%%%%%%%%%%%%%%%%%%%%%%%%%%%%%%%%%%%%%%%%%%%%%%%%%%%%%%%%%%%%%
% Lớp 10 | Phương pháp tọa độ trong mặt phẳng (𝑂𝑥𝑦) | Ba đường conic trong mặt phẳng tọa độ
% Dạng 0H9?5-7: Xác định các yếu tố của parabol
% Số câu: 7  (sinh tự động bởi tach_cau.py; từ nay thầy sửa trực tiếp file này)
%%%%%%%%%%%%%%%%%%%%%%%%%%%%%%%%%%%%%%%%%%%%%%%%%%%%%%%%%%%%%%%%%%%%%%

%%%%%%%%%%%%%%%% TRẮC NGHIỆM 4 PHƯƠNG ÁN %%%%%%%%%%%%%%%%

% nguồn: MaHoaNH_TracNghiem_Lop10
\begin{ex}%[0H9N5-7]%[ma:10230139]
	Trong mặt phẳng $Oxy$, cho phương trình parabol $(P)  \colon y^2=24x$. Tiêu điểm của parabol này là
	\choice 
		{$F(12;0)$} 
		{\True  $F(6;0)$}
		 {$F(-6;0)$}
		 {$F(-12;0)$}
	 \loigiai{
		Ta  có $2p=24\Leftrightarrow \dfrac{p}{2}=6$. Vậy tiêu điểm là $F(6;0)$.
	} 
\end{ex}

% nguồn: MaHoaNH_TracNghiem_Lop10
\begin{ex}%[0H9N5-7]%[ma:03f94409]
	Trong mặt phẳng tọa độ $Oxy$, cho parabol $(P)\colon y^2=16 x$. Khi đó tham  số tiêu của parabol $(P)$ bằng
	 \choice
		{\True  $8$}
		 {$4$}
		{$2$} 
		{$1$} 
	\loigiai{ 
		Ta có $y^2=16 x\Leftrightarrow y^2=2\cdot8 x$.\\ 
	Vậy parabol $(P)$ có  tham số tiêu $p=8$.}
\end{ex}


%%%%%%%%%%%%%%%% ĐÚNG/SAI %%%%%%%%%%%%%%%%

% nguồn: MaHoaNH_DungSai_Lop10
\begin{ex}%[0H9N5-7]%[ma:af13ea97]
	Cho parabol $(P)$ có phương trình chính tắc là $y^2=x$. Trong các mệnh đề dưới đây, mệnh  đề nào đúng, mệnh đề nào sai?
	 \choiceTF
		{\True Parabol có  tiêu điểm $F\left(\dfrac{1}{4};0\right)$}
		{\True Đường  chuẩn $\Delta \colon x=-\dfrac{1}{4}$}
		{Tham số tiêu của parabol  bằng $1$}
		{\True Các điểm có tọa độ $M_1\left(\dfrac{11}{4};\dfrac{\sqrt{11}}{2}\right)$; $M_2\left(\dfrac{11}{4};\dfrac{-\sqrt{11}}{2}\right)$ cách  đều điểm $F$}
	\loigiai{ 
		 Ta có $2p=1$. Do vậy
		\begin{itemchoice}
			 \itemch Tiêu điểm $F\left(\dfrac{1}{4};0\right)$.
			\itemch Đường  chuẩn $\Delta \colon x=-\dfrac{1}{4}$.
			\itemch Tham số tiêu $\mathrm{d}(F,\Delta)=p=\dfrac{1}{2}$. 
			\itemch Ta có  $\mathrm{d}(M_1,F)=p=\mathrm{d}(M_2,F)=3$.
		\end{itemchoice}
	 }
\end{ex}

% nguồn: MaHoaNH_DungSai_Lop10
\begin{ex}%[0H9N5-7]%[ma:6fe1a6a0]
	Cho parabol $(P)$ có phương trình chính tắc là $y^2=8x$.  Trong các mệnh đề dưới đây, mệnh đề nào đúng, mệnh đề nào sai?
	 \choiceTF
		{Parabol có  tiêu điểm $F\left(4;0\right)$}
		{\True Đường chuẩn $\Delta \colon  x=-2$}
		{Tham số  tiêu của parabol bằng $1$}
		{\True Các điểm $M_1\left(1;2\sqrt{2}\right)$; $M_2\left(1;-2\sqrt{2}\right)$ cách  đều $F$}
	 \loigiai{
		Ta có $2p=8$  nên $p=4$. Do vậy
		\begin{itemchoice}
			\itemch  Tiêu điểm $F\left(2;0\right)$.
			\itemch Đường chuẩn $\Delta \colon x=-2$. 
			\itemch Tham  số tiêu $\mathrm{d}(F,\Delta)=p=4$.
			\itemch Gọi điểm  $M_0(x_0; y_0)\in (P)$ sao cho $\mathrm{d}(M_0,F)=3=\mathrm{d}(M_0,\Delta)$. Hay
			$$\heva{&y_0^2=8x_0\\&\left|x_0+2\right|=3}\Leftrightarrow  \hoac{&x_0=1\Rightarrow y_0=2\sqrt{2}\\&x_0=1\Rightarrow y_0=-2\sqrt{2}.}$$
			Các điểm  $M_1\left(1;2\sqrt{2}\right)$; $M_2\left(1;-2\sqrt{2}\right)$ cách đều $F$.
		\end{itemchoice}
	} 
\end{ex}

% nguồn: MaHoaNH_DungSai_Lop10
\begin{ex}%[0H9V5-7]%[ma:a242f688]
	Trong mặt phẳng tọa độ $Oxy$, cho parabol $(P)\colon y^2=2x $ và hai điểm $ A(-3;2)$, $ B(0;-1) $. Gọi $C(a;b)$ là điểm thuộc parabol $ (P)$. Chọn khẳng định đúng, khẳng  định sai trong các khẳng định sau?
	\choiceTF 
		{\True Tọa độ của điểm  $C$ có dạng $ C\left(\dfrac{b^2}{2};b\right) $}
		 {\True Diện tích tam giác $ABC$ là $S=3\left(\dfrac{b^2}{4}+\dfrac{b}{2}+\dfrac{1}{2}\right)$}
		{\True Khi $ C\left(\dfrac{1}{2};-1\right)  $ thì tam giác $ABC$ có diện tích nhỏ nhất}
		{Khi tam giác $ABC$ cân tại $C$  thì diện tích tam giác $ABC$ đạt giá trị nhỏ nhất}
	 \loigiai{
		\begin{itemchoice}
			\itemch  Đây là khẳng định đúng vì\\
			Ta  có $ \overrightarrow{AB} =(3;-3)=3(1;-1)$.\\
			 Đường thẳng $ AB $ nhận $ \overrightarrow{u} =(1;1)$ làm một véc-tơ chỉ phương.\\
			 Suy ra đường thẳng $ AB $ nhận $ \overrightarrow{n} =(1;1)$ làm một véc-tơ pháp tuyến.\\
			Phương trình đường thẳng $ AB$ là $1\cdot (x-0)+1\cdot (y-(-1)) =0\Leftrightarrow  x+y+1=0$.\\
			Vì $ C(a;b) $ là  điểm thuộc parabol $ (P)$ nên $ b^2=2a \Rightarrow a=\dfrac{b^2}{2}$.\\
			Suy ra $  C\left(\dfrac{b^2}{2};b\right) $.\\
			 \itemch Đây là khẳng định đúng vì\\
			Diện tích tam giác $ ABC $ là 
			 \allowdisplaybreaks
			\begin{eqnarray*}
				S_{ABC}&=&\dfrac{1}{2}AB\cdot \mathrm{d}(C,AB)\\ 
				 &=&\dfrac{1}{2}\cdot \sqrt{3^2+(-3)^2}\cdot \dfrac{\left|\dfrac{b^2}{2}+b+1\right|}{\sqrt{1^2+1^2}}\\
				&=&\dfrac{3\sqrt{2}}{2}\cdot  \dfrac{\left|\dfrac{b^2}{2}+b+1\right|}{\sqrt{2}}\\
				 &=&\dfrac{3}{2}\left|\dfrac{b^2}{2}+b+1\right|=3\left(\dfrac{b^2}{4}+\dfrac{b}{2}+\dfrac{1}{2}\right).
			\end{eqnarray*}
			\itemch Đây là khẳng  định đúng vì \\
			Ta có 
			 $S_{ABC}=3\left(\dfrac{b^2}{4}+\dfrac{b}{2}+\dfrac{1}{2}\right)
			 =3\left(\left(\dfrac{b}{2}+\dfrac{1}{2}\right)^2+\dfrac{1}{4}\right)\geq 3\cdot \dfrac{1}{4}=\dfrac{3}{2}.$\\
			Ta có $ \min S_{ABC} =\dfrac{3}{2}$  đạt được khi và chỉ khi $b=-1$, $a=\dfrac{1}{2}$.\\
			Khi $  C\left(\dfrac{1}{2};-1\right) $ thì tam giác $ABC$ có diện tích nhỏ nhất.
			\itemch Đây là khẳng định sai  vì.\\
			Khi $ C\left(\dfrac{1}{2};-1\right) $ thì tam giác $ABC$  có diện tích nhỏ nhất.\\ Khi đó $CA=\dfrac{\sqrt{85}}{2}$, $CB=\dfrac{1}{2}$ nên tam giác $ABC$ không cân tại $C$.
		\end{itemchoice}
	 }
\end{ex}

% nguồn: MaHoaNH_DungSai_Lop10
\begin{ex}%[0H9V5-7]%[ma:73491d84]
	Trong mặt phẳng $Oxy$, cho parabol $(P)\colon y^2=8x$. Khi  đó, các khẳng định sau là đúng hay sai?
	 \choiceTF
		{\True Tiêu điểm  $F(2;0)$}
		 {\True Có hai điểm $M$ trên $(P)$, cách $F$ một khoảng là $3$}
		{Điểm $M$ trên $(P)$ sao cho $S_{\triangle OMF}=8$,  có hoành độ bằng $6$}
		{Tồn tại điểm $A$ nằm trên parabol và  điểm $B$ nằm trên đường thẳng $(\Delta)\colon 4x-3y+5=0$ sao cho đoạn thẳng $AB$ ngắn nhất, khi đó $AB$ ngắn nhất bằng $\dfrac{1}{5}$}
	\loigiai{ 
		\begin{center}
			\begin{tikzpicture}[scale=0.4,font=\footnotesize,line join=round,line cap=round,>=stealth]
					\def\xmin{-3}\def\xmax{5}\def\ymin{-4}\def\ymax{4}\pgfmathsetmacro\p{2} 
					\draw[->] (\xmin-0.2,0)--(\xmax+0.2,0) node[below] {\footnotesize  $x$};
					\draw[->]  (0,\ymin-0.2)--(0,\ymax+0.2) node[right] {$y$};
					\clip (\xmin,\ymin) rectangle  (\xmax,\ymax);
					\path 
						 (0,0) coordinate (O)
						 (-\p,0) coordinate (H) node[above left]{$\frac{-p}{2}$}
						(\p,0) coordinate (F) node[above  right]{$\frac{p}{2}$}
						(1,2) coordinate (M) 
						 (-\p,2) coordinate (N);
					 \draw
						(N)--(M)--(F) 
						 (-\p,\ymin)--(-\p,\ymax);
					\draw[smooth,variable=\x,blue] plot ({0.25*\x*\x},{\x}); 
					\fill  (M) node[right]{$M$} circle (1pt);
					 \foreach \x/\g in {H/-130,F/-90,O/-120}
					\fill (\x)  circle (1pt)
						($(\x)+(\g:5mm)$) node{$\x$}; 
				\end{tikzpicture}
		\end{center}
		\begin{itemchoice}
			\itemch \textbf{Đúng}.\\ Từ phương trình của $(P)$ suy ra $p=4$ nên tiêu điểm $F(2;0)$. 
			\itemch \textbf{Đúng}.\\ Giả sử $M(x_M; y_M)\in (P)$ suy ra  $y_M^2=8x_M$.\\
			Ta có $3=FM=x_M+\dfrac{p}{2}$ suy ra $x_M=1$,  thay vào $y_M^2=8x_M\Rightarrow y_M=\pm 2\sqrt{2}$.\\
			Vậy có hai  điểm thỏa mãn là $M_1(1;2\sqrt{2})$, $M_2(1;-2\sqrt{2})$.
			 \itemch \textbf{Sai}.\\ Ta có $M\in (P)$ nên $M\left(\dfrac{a^2}{8}; a\right)$ với $a\ge 0$. Suy ra $\mathrm{d}(M;OF)=y_M=a$.
			$$S_{\triangle OFM}=8\Leftrightarrow \dfrac{1}{2}OF\cdot  \mathrm{d}(M;OF)=8\Leftrightarrow a=8.$$
			Vậy điểm cần  tìm là $M(8;8)$.
			\itemch \textbf{Sai}.\\ Với mọi điểm $A\in (P)$, $B\in (\Delta)$, ta  luôn có $AB\ge \mathrm{d}(A;\Delta)$.\\
			$A\in (P)\Rightarrow A\left( \dfrac{a^2}{8};a\right)$ với $a\ge 0$. Khi đó,  $$\mathrm{d}(A;\Delta)=\dfrac{\left|4\cdot \frac{a^2}{8} -3\cdot a+5 \right|}{5}=\dfrac{(a-3)^2+1}{10}\ge \dfrac{1}{10}.$$
			Suy ra  $\mathrm{d}(A;\Delta)$ ngắn nhất khi $a=3$; $A\left( \dfrac{9}{8};3\right)$ và $B$ là hình chiếu của $A$ lên $(\Delta)$.\\
			Đường thẳng đi qua $A$ vuông góc với $(\Delta)$ nhận $\vec{u}=(3;4)$ làm véc-tơ pháp tuyến nên có  phương trình
			$$3\left(x-\dfrac{9}{8}\right)+4(y-3)=0\Leftrightarrow 24x+32y-123=0.$$ 
			Do đó tọa độ điểm $B$ là  nghiệm của hệ $\heva{& 4x-3y+5=0\\&24x+32y-123=0}\Leftrightarrow \heva{&x=\dfrac{209}{200}\\&y=\dfrac{153}{50}.}$\\
			Vậy có $A\left(  \dfrac{9}{8};3\right)$, $B\left( \dfrac{209}{200};\dfrac{153}{50}\right)$ để $AB$ ngắn nhất và $AB$ ngắn nhất bằng $\dfrac{1}{10}$.
		\end{itemchoice}
	 }
\end{ex}


%%%%%%%%%%%%%%%% TRẢ LỜI NGẮN %%%%%%%%%%%%%%%%

% nguồn: MaHoaNH_TraLoiNgan_Lop10
\begin{bt}%[0H9H5-7]%[ma:52669f87]
	Trong mặt phẳng tọa độ $Oxy$, cho parabol $(P)\colon y^2=2px$ ($p>0$). Biết khoảng  cách từ tiêu điểm $F$ đến đường thẳng $\Delta\colon x+y-12=0$ bằng $2\sqrt{2}$. Tính tổng các giá trị của $p$ thỏa mãn đề bài.
	\shortans{48} 
	\loigiai{Phương trình parabol $(P)\colon y^2=2px$ ($p>0$).  Tọa độ tiêu điểm của $(P)$ là $F\left( \dfrac{p}{2};0\right)$.\\
		 Khoảng cách từ $F$ đến đường thẳng $\Delta$ là $\mathrm{d}(F;\Delta)=2\sqrt{2}$.\\
		$$\mathrm{d}(F;\Delta)=2\sqrt{2}\Leftrightarrow \dfrac{\left|  \dfrac{p}{2}-12\right|}{\sqrt{2}}=2\sqrt{2}\Leftrightarrow \left| \dfrac{p}{2}-12\right|=4\Leftrightarrow \heva{&\dfrac{p}{2}-12=4\\&\dfrac{p}{2}-12=-4}\Leftrightarrow \hoac{&p=32\\&p=16.}$$
		Vậy Tính tổng các giá trị của $p$ thỏa mãn  đề bài là $T=32+16=48$.
	} 
\end{bt}
